Let the turning points be $U,V$. The points $P,U,V$ lie on the circle centered at $Q$ with radius $10$. Let $R_0$ be the point diametrically opposite $P$, namely $(0,20)$. By Thales' theorem, $UR_0\perp UP$ and $VR_0\perp VP$. The designated turns point toward $R_0$: the direction from $U$ is $65^\circ$, and from $V$ it is $145^\circ$. Thus the two outgoing rays meet at $R=R_0$, giving $\boxed{PR=20}$.
\begin{asy}
texpreamble("\usepackage{amsmath,amssymb}");
size(240);
defaultpen(fontsize(9pt));
pair P=(0,0),Q=(0,10),R=(0,20),U=20*Sin(25)*dir(155),V=20*Sin(55)*dir(55);draw(circle(Q,10),gray);draw(P--U--R,blue,Arrows);draw(P--V--R,red,Arrows);draw(P--R,dashed);label("$P$",P,(0,-1));label("$Q$",Q,(-1,0));label("$R$",R,(0,1));label("$U$",U,(-1,0));label("$V$",V,(1,0));
\end{asy}
